\section{Manifolds with singularities}\label{sec6}
All results on the heat kernel expansion formulated in the
previous sections are valid on
smooth manifolds only. If there are boundaries, they also
have to be smooth. As well, any singularities in the potential
term or in the field strength are strictly speaking forbidden.
However, many physical models deal with singular backgrounds.
Even if such backgrounds may be represented through certain limiting
procedures from smooth configurations, the heat kernel coefficients
are not given by limits of their ``smooth'' values. The most visible
manifestation of failure of the smooth field approximation is that $a_k$
with sufficiently large $k$ are divergent. Usually, the presence of
singularities changes even the structure of the heat kernel
kernel expansion as compared to the smooth case. 

%%%%%%%%%%
\subsection{Non-integrable potentials}\label{s6sinpot}
According to (\ref{divWLam}) divergences in the effective action
are defined by {\it integrated} heat kernel coefficients.
Although the formulae (\ref{a0nobou}) - (\ref{a6nobou}) for 
the {\it localised} heat kernel coefficients are valid on
non-compact manifolds (provided the smearing function $f$
falls off sufficiently fast), transition to the integrated
heat kernel is not that straightforward. Already the coefficient
$a_0(D)=a_0(1,D)$, which is proportional to the volume,
is divergent. This divergence is usually removed by replacing
$\det (D)$ in (\ref{pathdet}) by
\begin{equation}
\det (D)/\det (D_0) \,,\label{referdet}
\end{equation}
where the operator $D_0=-\partial^2+m^2$ describes a free particle
propagation in an ``empty'' space. It is argued that since $D_0$ does
not depend on ``essential'' variables division by $\det (D_0)$ does not
change physical predictions of the theory. In all subsequent formulae
the heat kernel $K(t;x,y;D)$ is then replaced by the subtracted
heat kernel
\begin{equation}
K_{\rm sub} (t;x,y)=K(t;x,y;D)-K(t;x,y;D_0) \,.
\label{Ksub}
\end{equation}
In flat space the coefficient $a_0$ corresponding to $K_{\rm sub}$
is identically zero\footnote{On a curved manifold the subtraction
procedure is more subtle. On has to define a reference metric which
differs from the physical one on a compact submanifold.}. If the field
strength $\Omega_{\mu\nu}$ and the (subtracted) potential $E+m^2$ 
have a compact support or decay sufficiently fast at the infinity,
the small $t$ asymptotic expansion of $K_{\rm sub}(t;x,x)$ is
integrable on the whole $M$. If not, the very structure of the
global heat kernel may be changed.

As an example of non-integrable potentials consider the harmonic
oscillator in one dimension. The Schr\"{o}dinger operator reads
\begin{equation}
D=-\partial_x^2 +\nu^2 x^2 \,.\label{harmos}
\end{equation}
If we consider the problem on the whole real axis, $M=\mathbb{R}$,
the potential term is not integrable. Already the expression
(\ref{a2nobou}) for $a_2(1,D)$ diverges. Therefore, analytic expressions
of sec.\ \ref{sec4} cannot be used in this case. However, the
(integrated) heat kernel can be easily calculated. Eigenvalues of the
operator (\ref{harmos}) are contained in almost any textbook on quantum
mechanics:
\begin{equation}
\lambda_j=\nu (2j+1),\qquad j=0,1,2,\dots \label{harmspec}
\end{equation}

The integrated heat kernel reads
\begin{equation}
K(t;D)=\sum_{j=0}^\infty e^{-t\nu (2j+1)} =
\frac 12 [\sinh (\nu t)]^{-1} \,.\label{harmheat}
\end{equation}
As $t\to 0$ it behaves like $1/t$ while for smooth rapidly decaying
potentials in one dimension the leading singularity in the heat kernel
is $1/\sqrt{t}$. This statement may be generalised to higher dimensions.
If $D=-\partial^2 +P_{\mu\nu}x^\mu x^\nu$ with a non-degenerate 
matrix $P_{\mu\nu}$ on $M=\mathbb{R}^n$, the leading
term in $K(t;D)$ is $(2t)^{-n} (\det P)^{-1/2}$ \cite{Avramidi:1995ik}.

\subsection{Conical singularities}\label{s6consin}
Conical space is defined as $M=[0,1]\times N$ where $N$ is an
$n-1$-dimensional manifold called the base. The metric of the
cone has the form
\begin{equation}
(ds)^2=dr^2 +r^2 d\Omega^2 \,,\label{conmetr}
\end{equation}
where $r\in [0,1]$ and $d\Omega^2$ is the line element on the
base $N$. This metric is, in general, singular at $r=0$. However, if we 
take the unit sphere $S^{n-1}=N$ with standard round metric,
the singularity disappears and we obtain the $n$-dimensional unit ball
(\ref{ballmetric}). If a manifold has singular points where
the metric can be approximated by (\ref{conmetr}) we say that this
manifold has conical singularities. 

Conical singularities appear in many physical applications.
First of all, with $N=S^1$ the metric (\ref{conmetr}) is the Euclidean
version of the Rindler metric. Conical singularities appear in
classical solutions of the Einstein equations 
\cite{Sokolov:1977,Frolov:1987dz,Barriola:1989hx}
and in the supermembrane theory 
\cite{Vassilevich:1991vy}. Gravitational field of a point mass in 
three dimensional gravity is a conical space 
\cite{Deser:1984tn,Deser:1988qn}. There are evidences \cite{Schleich:1993bs}
that ``conifolds'' dominate the path integral for quantum gravity in
topological sectors.

Sommerfeld \cite{Sommerfeld:1894}
was probably the first to consider the heat kernel in the presence
of conical singularities. The mathematical theory of the heat kernel
asymptotics with conical singularities was developed almost 100 years
later \cite{Cheeger:1983,Bruening:1984,Bruening:1987}. There two peculiar
features of these asymptotics. First, the heat kernel expansion
contains in general both integer and half-integer powers of $t$ even
without boundaries. Second, a non-standard $\ln t$ term may be contained in
the asymptotic series\footnote{The present author is not aware of
any simple example where the $\ln t$ terms actually appear.}.

On a manifold with conical singularities no closed analytical expression
for the heat kernel coefficients is available. However, usually it is
possible to disentangle contributions of the singularities from the
smooth part. For example, if $N=S^1$, $d\Omega^2=d\varphi^2$ with
$\varphi \in [0,\alpha ]$ only $a_2$ receives a contribution from the
tip of the cone:
\begin{equation}
a_2({\rm tip})=\frac{4\pi^2 -\alpha^2}{24\pi \alpha} \,.\label{a2tip}
\end{equation} 

In many particular cases of conical singularities
a very detailed analysis of the heat kernel
expansion has been performed 
\cite{Dowker:1977zj,Dowker:1989pp,Chang:1993fu,Dowker:1994jv,
Fursaev:1994qk,Fursaev:1994in,Cognola:1994qg,Cognola:1994ps,
Fursaev:1995ef,Bordag:1996fw,DeNardo:1997kp,Shtykov:1995fe,
Dowker:1998sc,BezerradeMello:2000mb}. One-loop computations on
general orbifolds were considered recently in 
\cite{GrootNibbelink:2003gd}.
%%%%%%%
\subsection{Domain walls and brane world}\label{s6dom}
Delta function is an example of an extremely sharp background potential.
Let us consider a manifold $M$ and a submanifold $\Sigma$ of the
dimension $n-1$. Let
\begin{equation}
D[\dvxi ]=D+\dvxi \delta_\Sigma \,.\label{singop}
\end{equation}
$D$ is an operator of Laplace type (\ref{laplaceb}). 
Let $h$ be
the determinant of the induced metric on $\Sigma$. Then $\delta_\Sigma$ is a
delta function defined such that
\begin{equation}
\int_M dx\sqrt{g} \delta_\Sigma f(x) =\int_\Sigma dx \sqrt{h}
f(x) \,.\label{deltaSig}
\end{equation}

The spectral problem
for $D[\dvxi ]$ on $M$ as it stands is ill-defined 
owing to the discontinuities (or
singularities) on $\Sigma$. It should be replaced by a pair of spectral
problems on the two sides
$M^\pm$ of
$\Sigma$ together with suitable matching conditions on $\Sigma$. In order to
find such matching conditions, we 
consider an eigenfunction $\phi_\lambda$ of the operator 
(\ref{singop}):
\begin{equation}
D[v]\phi_\lambda =\lambda \phi_\lambda \,.
\label{eigen}
\end{equation}
Let us choose the coordinate system such that $e_n$ is a unit normal
to $\Sigma$ and $x^n=0$ on $\Sigma$.
It is clear that $\phi_\lambda$ must be continuous on $\Sigma$:
\begin{equation}
\phi\vert_{x^n=+0}=\phi\vert_{x^n=-0} \,.\label{mc1}
\end{equation}
Otherwise, the second normal derivative of $\phi_\lambda$ would create
a $\delta'$ singularity on $\Sigma$ which is absent on the right
hand side of (\ref{eigen}). Let us integrate (\ref{eigen}) over
a small cylinder $\mathcal{C}=C^{n-1}\times [-\epsilon ,+\epsilon]$
\begin{equation}
\int_{\mathcal{C}} d^nx\sqrt{g} \left( -\nabla_n^2 \phi_\lambda
-\left[ \nabla_a^2 \phi_\lambda +(E+\lambda)\phi_\lambda \right]\right)
+\int_C d^{n-1}x\sqrt h \dvxi\phi_\lambda =0\,. \label{intcyl}
\end{equation} 
We now take the limit as $\epsilon\to 0$. Since the expression in the
square brackets in (\ref{intcyl}) is bounded, the contribution
that this term makes vanishes in the limit. We obtain
\begin{equation}
0=\int_C d^{n-1}x\sqrt h \left( -\nabla_n \phi_\lambda\vert_{x^n=+0}
+\nabla_n \phi_\lambda\vert_{x^n=-0} +\dvxi\phi_\lambda \right)\,.
\label{int2cyl}
\end{equation}
Since $C$ and $\lambda$ are arbitrary, we conclude that a proper
matching condition for the normal derivatives is
\begin{equation}
-\nabla_n \phi \vert_{x^n=+0}
+\nabla_n \phi \vert_{x^n=-0} +\dvxi\phi =0 \,.\label{dermatch}
\end{equation}

Physically this problem corresponds to two domains separated by a
penetrable membrane $\Sigma$ (a domain wall). In many cases penetrable
membranes are better models of physical boundaries then just boundary
conditions which are imposed on each side of $\Sigma$ independently
and thus exclude any interaction between the domains
\cite{Jaekel:1992bn,Actor:1995vc}. $\delta$-potentials
are being used in quantum mechanics \cite{Albeverio:1988}
where one studies the Schr\"{o}dinger equation (which is
nothing else than the imaginary time heat equation).
The Casimir energy calculations have been performed e.g. in
\cite{Scandurra:1998xa}.
In the formal limit
$\dvxi\to\infty$ one obtains Dirichlet boundary conditions 
on $\Sigma$, although the heat kernel coefficients are divergent in this
limit (see below).

Further generalisations are suggested by the brane world scenario
\cite{Randall:1999ee,Randall:1999vf} which assumes that our world
is a four dimensional membrane in a five dimensional space\footnote{
A similar scenario was proposed earlier in \cite{Rubakov:1983bb},
see \cite{Rubakov:2001kp} for a review.}. According to the Israel
junction condition \cite{Israel:1966rt} the metric in such models
cannot be smooth on $\Sigma$. Typical form of the metric near 
$\Sigma$ is 
\begin{equation}
(ds)^2=(dx^n)^2 + e^{-\alpha |x^n|} (ds_{n-1})^2, \label{branemetr}
\end{equation}
where $\alpha$ is a constant and where $(ds_{n-1})$ is a line element
on the $(n-1)$-dimensional hypersurface $\Sigma$. Due to the presence of the
absolute value of the $n$-th coordinate in (\ref{branemetr}), the normal
derivative of the metric jumps on $\Sigma$.
One can think of two smooth manifolds
$M^+$ and $M^-$ glued together along their common boundary $\Sigma$.
Neither Riemann tensor, nor matrix potential $E$ must be continuous
on $\Sigma$. Also, the extrinsic curvatures $L_{ab}^+$ and $L_{ab}^-$
of $\Sigma$ considered as a submanifold in $M^+$ and in $M^-$ 
respectively are, in general,
different. All geometric quantities referring to $M^-$ (respectively,
$M^-$) and their limiting values on $\Sigma$ will be supplied by
a superscript ``$+$'' (respectively, ``$-$'').  

For the case at hand there is still an asymptotic expansion (\ref{asymptotex})
for the heat kernel. The heat kernel coefficients can be decomposed as
\begin{equation}
a_k(f,D[\dvxi ])=a_k^+(f,D) + a_k^-(f,D)^- +a_k^\Sigma (f,D,\dvxi )\,,
\label{apmSigma}
\end{equation}
where $a_k^\pm (f,D)$ are known volume contributions corresponding
to $M^\pm$ (cf. (\ref{a0nobou}) - (\ref{a6nobou})). The coefficients
$a_k^\Sigma$ are given by integrals over $\Sigma$ of some local
invariants. Note, that 
\begin{equation}
\tilde\omega_a=\nabla_a^+-\nabla_a^- \label{tiloma}
\end{equation}
being a difference of two connection is a (pseudo-) vector with respect to
all space-time and gauge symmetries. Therefore, $\tilde\omega_a$ can be
used for constructing the surface invariants.

To make the formulae more symmetric we introduce {\it two} inward pointing
unit normals $\nu^+$ and $\nu^-$ to $\Sigma$ in $M^+$ and $M^-$ respectively.
We do not suppose that the smearing function $f$ is smooth on $\Sigma$
(therefore, there is no relation between $f_{;\nu^+}$ and $f_{;\nu^-}$),
but we assume continuity of $f$: $f^+=f^-=f$ on $\Sigma$.  

The surface invariants can be constructed from $L_{ab}^\pm$,
$R_{ijkl}^\pm$, $E^\pm$, $\dvxi$, $\tilde\omega_a$ and their
derivatives. This gives much more invariants than we usually have
for a boundary value problem. There are, however, some properties
of $a_k^\Sigma$ which simplify the calculations considerably.
First of all, $a_k^\Sigma$ must be invariant with respect to
interchanging the roles of $M^+$ and $M^-$. Also, $a_k^\Sigma$
must vanish when the singularity disappears. The first property
excludes, for example, the term $f(E^+-E^-)$ which changes sign
under $M^+\leftrightarrow M^-$. The second requirement excludes
the invariant $f(E^++E^-)$ because it survives even if there is no
singularity on $\Sigma$. These simple arguments show that $a_3^\Sigma$
does not contain $E^\pm$ even though such terms are allowed on dimensional
grounds. It is also very helpful that in some particular cases the
problem in question can be reduces to a sum of Dirichlet and Robin
boundary value problems \cite{Bordag:1999ed,Gilkey:2001mj}.

The coefficients $a_k^\Sigma$, $k=0,1,2,3$, read
\begin{eqnarray}
&&a_0^\Sigma(f,D,\dvxi)=0.\nonumber\\
&&a_1^\Sigma(f,D,\dvxi)=0.\nonumber\\
&&a_2^\Sigma(f,D,\dvxi)=
     (4\pi)^{-n/2}\frac16\int_\Sigma d^{n-1}x\sqrt{h}
       \ptr \{2f(L_{aa}^++L_{aa}^-)-6f\dvxi\}.\label{asig03}\\
&&a_3^\Sigma(f,D,\dvxi)=(4\pi)^{(1-n)/2}\frac1{384}
\int_\Sigma d^{n-1}x\sqrt{h}
\ptr\left\{\frac32f(L_{aa}^+L_{bb}^+
+L_{aa}^-L_{bb}^-\right.  \nonumber\\
&&\qquad +2L_{aa}^+L_{bb}^-)
    +3f(L_{ab}^+L_{ab}^++L_{ab}^-L_{ab}^-+2L_{ab}^+L_{ab}^-)
\nonumber\\
&&\qquad    +9(L_{aa}^++L_{aa}^-)(f_{;\nu^+}^++f_{;\nu^-}^-)
+48f\dvxi^2+24f\tilde\omega_a\tilde\omega_a \nonumber\\
&&\qquad \left.   -24f(L_{aa}^++L_{aa}^-)\dvxi
    -24(f_{;\nu^+}^++f_{;\nu^-}^-)\dvxi\right\}.\nonumber
\end{eqnarray}
The coefficients $a_4^\Sigma$ and $a_5^\Sigma$ are too long to be
presented here in full generality. Therefore, we restrict ourselves to 
the case of smooth geometry ($R_{ijkl}^+=R^-_{ijkl}$, $L_{ab}^+=-L_{ab}^-$),
smooth connection ($\tilde\omega_a=0$), and smooth smearing function
($f_{;\nu^+}=-f_{;\nu^-}=f_{;n}$). In other words, the only singularity comes
from the surface potential $\dvxi$.
\begin{eqnarray}
&&a_4^\Sigma (f,D,\dvxi)=(4\pi )^{-n/2}
\int_\Sigma d^{n-1}x\sqrt{h}\ptr\left\{ -\frac 16 f\dvxi^3 -\frac 1{6}
fR \dvxi -fE\dvxi \right. \nonumber \\
&&\qquad \left. -\frac 16 f\dvxi_{:aa} +\frac 16 f_{;n}\dvxi L_{aa} 
-\frac 16 f_{;nn}\dvxi
\right\} \nonumber \\
&&a_5^\Sigma (f,D)= (4\pi )^{-(n-1)/2} \int_\Sigma d^{n-1}x\sqrt{h}\ptr
\left\{\frac 1{64} f\dvxi^4 +\frac 1{48} fR \dvxi^2
 \right. \label{asig45}\\
&&\qquad +\frac 1{192}R_{nn}\dvxi^2 +\frac 18 f\dvxi^2 E 
-\frac 1{256} f\dvxi^2 L_{aa}L_{bb} 
+\frac 1{128}f\dvxi^2 L_{ab}L_{ab}  \nonumber \\
&&\qquad \left. +\frac 1{24}f\dvxi_{:aa}\dvxi
+\frac 5{192} f\dvxi_{:a}\dvxi_{:a} -\frac 5{384}f_{;n} \dvxi^2L_{aa}
+\frac 1{64}f_{;nn}\dvxi^2 \right\}
\nonumber
\end{eqnarray}
The heat kernel coefficients for $\Sigma =S^{n-1}\subset\mathbb{R}^n$
and $\dvxi =const.$ were calculated in \cite{Bordag:1998vs}.
Generic $\Sigma$ with arbitrary $\dvxi$ was considered in
\cite{Bordag:1999ed}. Moss \cite{Moss:2000gv} added a non-smooth
connection (see also \cite{Drozdov:2002um}).
Calculations on a particular brane-world background can be found
in \cite{Nojiri:2000bz}.
The heat kernel coefficients in the general setting
described here were calculated in \cite{Gilkey:2001mj}. This latter
paper also considered renormalization of the brane-world scenario
and predicted a non-standard Higgs potential on the brane.
Related calculations in 
wormhole models were done in 
\cite{Khusnutdinov:2002qb,Khusnutdinov:2003ii}. The $\zeta$
function for brane-world geometries with matching conditions
(\ref{mc1}), (\ref{dermatch}) was considered recently in
\cite{Elizalde:2002dd,Moss:2003zk} 
(where one can find some further references).

It is very well known \cite{Albeverio:1988}
that the conditions (\ref{mc1}),
(\ref{dermatch}) (which we call {\it transmittal}) are not the most
general matching conditions which can be defined on a surface. 
In general, boundary values of a function and of its' normal
derivatives are related by a $2\times 2$ {\it transfer} matrix:
\begin{equation}
0=
     \left(\begin{array}{cc}
          \nabla_{\nu^+}^++S^{++},\qquad&S^{+-}\\
          S^{-+},\qquad&\nabla_{\nu^-}^-+S^{--}\end{array}\right)
      \left(\begin{array}{l}\phi^+\\\phi^-\end{array}\right)
\bigg|_\Sigma\,.\label{BBT}
\end{equation} 
Note, that the transfer conditions (\ref{BBT}) do not assume
identification of $\phi^+$ and $\phi^-$ on $\Sigma$. In other
words, there is no ad hoc relation between the restrictions
of the vector bundles $V^+|_\Sigma$ and $V^-|_\Sigma$. We can even
consider the situation when we have ${\rm dim}V^+\ne{\rm dim}V^-$,
i.e. the fields on $M^-$ and $M^-$ can have different structures
with respect to space-time and internal symmetries. $S^{\pm\pm}$
are some matrix valued functions on $\Sigma$ (one can even consider
the case when they are differential operators). It is clear from
the notations on which spaces they act. For example,
$S^{+-}: V^-|_\Sigma\to V^+|_\Sigma$.

The matching conditions (\ref{BBT}) arise in heat transfer problems
\cite{Carslaw:1986}, some problems of quantum mechanics 
\cite{Grosche:1994uv}, and in conformal field theory
\cite{Bachas:2001vj}. In a formal limiting case 
$S^{++}-S^{-+}=S^{--}-S^{+-} \to \infty$ while 
$\dvxi =2(S^{++}+S^{+-})$ is kept finite one arrives at transmittal
boundary conditions (\ref{mc1}), (\ref{dermatch}). The heat kernel
coefficients are divergent in this limit.

In a particular case of spherical $\Sigma$ the heat kernel
expansion with transfer boundary conditions was evaluated in 
\cite{Bordag:2001ta}. General expressions for $a_k$ with
$k=0,1,2,3,4$ were obtained in \cite{Gilkey:2002nv}.
Somewhat surprisingly, the calculations for (\ref{BBT})
are easier than for the singular particular case (\ref{mc1}),
(\ref{dermatch}).

We should note that not all singular limiting cases of
(\ref{BBT}) are described by the transmittal conditions
(\ref{mc1}), (\ref{dermatch}). The heat kernel expansion
for a generalisation of transmittal condition is known
in the spherically symmetric case only \cite{Bordag:2001ta}.  
Very little is known about the heat kernel if the transfer
matrix contains differential operators on $\Sigma$ (conformal
walls of ref.\ \cite{Bachas:2001vj} belong to this class
of problems).

The case when the singular potential is located on a surface
of co-dimension larger than one (i.e. when 
${\rm dim}\Sigma <n-1$) is rather complicated. Even a careful
translation of this problem to the operator language
was done only in 1960's \cite{Berezin:1961}. Direct calculations
show \cite{Albeverio:1994,Solodukhin:1997xn} that
the heat kernel asymptotics may contain very unusual $(\ln t)^{-1}$
terms. More references can be found in \cite{Camblong:2001um}.
%%%%%%
\subsection{Non-smooth boundaries}\label{s6nsbou}
Rectangular region in a plane is probably the simplest
manifold with boundaries as far as eigenvalues of the Laplacian  
are concerned. However, the formulae (\ref{a0bou}) - 
(\ref{a4bou}) are not valid for this case because of the presence
of corners. The heat kernel expansion on manifolds
with piecewise smooth boundaries was considered by Kac in his
famous paper \cite{Kac:1966}. He demonstrated that for a 
region in $\mathbb{R}^2$ each corner with the inside
facing angle $\alpha$ contributes 
\begin{equation}
a_2({\rm corner})= \frac {\pi^2 -\alpha^2}{24\pi\alpha}
\label{a2Kac}
\end{equation}
to the coefficient $a_2$ while $a_0$ and $a_1$ are still defined
by their ``smooth'' expressions.

The formula (\ref{a2Kac}) looks similar to the contribution of a
conical singularity (\ref{a2tip}). The reason for this similarity
is that the cone can be obtained from the wedge by gluing the sides
together and imposing the periodicity conditions.

The study of boundary discontinuities was continued by Apps and
Dowker who calculated the coefficients $a_3$ \cite{Dowker:1995sp}
and $a_4$ \cite{Apps:1998zr} for piecewise smooth boundaries.
We also refer to \cite{Aurell:1994} where functional determinants
on simplicial complexes were analysed, and to
\cite{Dowker:2000fs} where divergences in the Casimir energy
found in \cite{Nesterenko:2000eb} were attributed to non-smoothness
of the boundaries. A recent study \cite{Nesterenko:2002ng}
should also be mentioned.

When the angle $\alpha$ goes to $0$ we obtain a cusp. In this limit
(\ref{a2Kac}) is divergent. Presence of the cusp is an essential
singularity which modifies {\it powers} of the proper
time $t$ which appear in the asymptotic expansion of the heat
kernel \cite{Stewartson:1971}.
%%%%%%%%%
\subsection{Dielectric bodies}\label{s6diel}
Calculations of the Casimir energy of a dielectric body 
have attracted much attention and created many controversial
results. A (rather large) literature on this subject is reviewed
in \cite{Bordag:2001qi,Milton:2001yy}. Quantum field theory formulation
of this problem is known for a long time already (cf. 
\cite{Schwinger:1979pa,Milton:1980yx,Candelas:1982qw}). However, the
heat kernel analysis of divergences in the Casimir energy in
dielectric is a relatively new subject. 

Wave propagation with variable speed of light $c(x)$
is described by the operator $D=-c(x)^2 \nabla^2 +\dots$,
where we neglected the lower order terms. In a dielectric medium
$c$ is expressed in terms of the dielectric permittivity $\varepsilon$
and of the magnetic permeability $\mu$:
$c(x)^2=1/(\varepsilon (x)\mu (x))$. 
For a smooth distribution of $\varepsilon$ and $\mu$
the operator $D$
is a particular case of curved space Laplacian with
an effective metric defined by $c(x)$. The heat kernel
coefficients can be calculated in the standard way \cite{Bordag:1998fh}.

Consider now a dielectric body bounded by $\Sigma$. Typically, 
$c(x)$ (and the effective metric!) jumps on $\Sigma$.
This singularity is much stronger than the one considered
in sec.\ \ref{s6dom}. Thus the geometric interpretation of this
problem is very difficult. Very little is known about the heat kernel
expansion in a dielectric body of an arbitrary shape.
There are calculations for a dielectric ball 
\cite{Bordag:1998vs} and for a dielectric cylinder \cite{Bordag:2001zj}.
These calculations exhibit a puzzling property of the heat kernel
expansion in dielectrics in dilute approximation ($\varepsilon \sim 1$):
for a dielectric body the Casimir energy
in the ultra violet limit behaves 
better than for the ``smooth'' case.
The heat kernel expansion for a frequency dependent $\varepsilon$
was considered in refs.\ \cite{Falomir:2001uv,Bordag:2001fq}.
